\subsection{Seleção de Repositórios}

Foram selecionados 197 repositórios populares do GitHub, extraídos da busca
pública ordenada por estrelas. Critérios de inclusão: pelo menos 100 PRs fechados.

\subsection{Seleção de Pull Requests}

Foram analisados 9.075 PRs que atendem aos seguintes critérios:
\begin{itemize}
    \item Status MERGED ou CLOSED;
    \item Pelo menos uma revisão oficial (campo \texttt{review\_count} $\geq$ 1);
    \item Tempo de análise maior que 1 hora.
\end{itemize}

\subsection{Métricas}

As métricas extraídas para cada PR são:
\begin{itemize}
    \item \textbf{Tamanho:} número de arquivos modificados e total de linhas alteradas;
    \item \textbf{Tempo de Análise:} intervalo entre criação e fechamento/merge (horas);
    \item \textbf{Descrição:} número de caracteres do corpo do PR;
    \item \textbf{Interações:} número de participantes distintos;
    \item \textbf{Revisões:} número de submissões de revisão oficiais.
\end{itemize}

\textbf{Nota:} os campos \texttt{comments\_count} e \texttt{review\_comments\_count}
apresentaram valores zerados na coleta. Como proxy para interações,
utilizamos o número de participantes, que foi corretamente coletado.

\subsection{Testes Estatísticos}

Utilizamos o coeficiente de correlação de \spearman ($\rho$), um método
não-paramétrico robusto a outliers e distribuições não-normais.
Aplicamos winsorização no percentil 99 para mitigar o impacto de outliers extremos.

O nível de significância adotado foi $\alpha = 0.05$. A magnitude do efeito
foi classificada conforme o valor de $|\rho|$:
desprezível ($<$0,1), fraca (0,1--0,3), moderada (0,3--0,5), forte (0,5--0,7),
muito forte ($>$0,7).

\subsection{Visualizações}

Para a representação gráfica, seguimos as recomendações de Claus O. Wilke
(\textit{Fundamentals of Data Visualization}, O'Reilly, 2019):
\begin{itemize}
    \item \textbf{ECDF} (Função de Distribuição Cumulativa Empírica) para comparações
    de distribuição entre grupos (Dimensão A), por ser livre de parâmetros arbitrários
    (como largura de bins em histogramas) e robusta a outliers.
    \item \textbf{Violin plots com boxplot interno} para relações entre variável contínua
    e discreta (Dimensão B), revelando a distribuição completa sem adicionar ruído artificial.
    \item \textbf{Matriz de correlação triangular} para visão geral das relações
    entre todas as métricas.
\end{itemize}
